\documentclass[aspectratio=169]{beamer}
\usetheme{metropolis} % modern dark-friendly theme
\usepackage[utf8]{inputenc}
\usepackage[serbian]{babel}
\usepackage{amsmath, amssymb}
\usepackage{listings}
\usepackage{xcolor}

% Dark theme adjustments
\definecolor{bgcolor}{RGB}{20,20,25}
\definecolor{codegray}{gray}{0.9}
\definecolor{mintgreen}{RGB}{0,255,180}
\setbeamercolor{background canvas}{bg=bgcolor}
\setbeamercolor{normal text}{fg=white}

\lstset{
  backgroundcolor=\color{black},
  basicstyle=\ttfamily\footnotesize\color{white},
  keywordstyle=\color{mintgreen},
  commentstyle=\color{gray},
  stringstyle=\color{orange},
  showstringspaces=false,
  frame=single,
  rulecolor=\color{gray},
  breaklines=true
}

\title{Kvantno računarstvo}
\author{Stefan Nožinić}
\date{Oktobar 2025}

\begin{document}

\maketitle

% --- Slide 1 ---
\begin{frame}{Klase problema u računarstvu}
\begin{itemize}
    \item Pojmovi: P, NP, NP-težak, NP-kompletan, EXP
    \item Zašto su ovi pojmovi važni?
\end{itemize}
\end{frame}

% --- Slide 2 ---
\begin{frame}{Klasa P}
\begin{itemize}
    \item P (Polynomial time) je klasa problema koji se mogu rešiti u polinomijalnom vremenu na determinističkom Turingovom mašinom.
    \item Primeri problema u P:
    \begin{itemize}
        \item Sortiranje niza brojeva
        \item Pronalaženje najkraćeg puta u grafu (Dijkstra algoritam)
        \item Provera da li je broj paran
        \item Množenje dva broja
        \item Pretraga u sortiranoj listi (binarna pretraga)
        \item Računanje najvećeg zajedničkog delioca (GCD)
    \end{itemize}
    \item Problemi u P su efikasno rešivi na klasičnim računarima.
\end{itemize}
\end{frame}

% --- Slide 3 ---
\begin{frame}{Klasa NP}
\begin{itemize}
    \item NP (Nondeterministic Polynomial time) — problemi čija se rešenja mogu proveriti u polinomijalnom vremenu.
    \item Primeri:
    \begin{itemize}
        \item Svi problemi iz klase P
        \item TSP
        \item SAT
        \item Problem bojenja grafa
        \item Faktorizacija velikih brojeva
    \end{itemize}
    \item Otvoreno pitanje: Da li je P = NP?
\end{itemize}
\end{frame}

% --- Slide 4 ---
\begin{frame}{NP-težak i NP-kompletan}
\begin{itemize}
    \item NP-težak (NP-hard): bar toliko težak kao i najteži problemi u NP.
    \item NP-kompletan (NP-complete): u NP i NP-težak.
    \item Ako se bilo koji NP-kompletan problem reši u polinomijalnom vremenu, P = NP.
    \item Primeri:
    \begin{itemize}
        \item SAT
        \item Problem bojenja grafa
        \item TSP
    \end{itemize}
\end{itemize}
\end{frame}

% --- Slide 5 ---
\begin{frame}{Klasa EXP}
\begin{itemize}
    \item EXP — problemi rešivi u eksponencijalnom vremenu.
    \item Primeri:
    \begin{itemize}
        \item Problemi kombinatorike i optimizacije
        \item Problemi koji zahtevaju ispitivanje svih kombinacija
    \end{itemize}
    \item Veza: svi problemi u NP su u EXP, ali obrnuto nije poznato.
\end{itemize}
\end{frame}

% --- Slide 6 ---
\begin{frame}{Kvantno računarstvo i klase problema}
\begin{itemize}
    \item Kvantni računari koriste kubite i fenomene superpozicije i uplitanja.
    \item Mogu rešavati određene probleme efikasnije:
    \begin{itemize}
        \item Šorov algoritam — faktorizacija
        \item Groverov algoritam — pretraga
    \end{itemize}
\end{itemize}
\end{frame}

% --- Slide 7 ---
\begin{frame}{Double Slit Experiment}
\begin{itemize}
    \item Pokazuje dualnost talasa i čestica.
    \item Osnova za razumevanje superpozicije i kvantnih stanja.
\end{itemize}
\end{frame}

% --- Slide 8 ---
\begin{frame}{Kubiti i superpozicija}
\begin{itemize}
    \item Kubit: osnovna jedinica kvantne informacije.
    \item Može biti u stanju $|0\rangle$, $|1\rangle$ ili superpoziciji oba:
\end{itemize}

\[
|\psi\rangle = \alpha|0\rangle + \beta|1\rangle
\]

\[
|\alpha|^2 + |\beta|^2 = 1
\]
\end{frame}

% --- Slide 9 ---
\begin{frame}{Inherentni paralelizam kvantnog računarstva}
\begin{itemize}
    \item Sa $n$ kubita, moguće je istovremeno obrađivati $2^n$ stanja.
    \item Kvantni paralelizam omogućava simultano izračunavanje funkcije $f(x)$ za sve ulaze.
\end{itemize}
\end{frame}

% --- Slide 10 ---
\begin{frame}{Ket notacije}
\[
|0\rangle = \begin{pmatrix} 1 \\ 0 \end{pmatrix}, \quad
|1\rangle = \begin{pmatrix} 0 \\ 1 \end{pmatrix}
\]
\[
|\psi\rangle = \alpha \begin{pmatrix}1 \\ 0\end{pmatrix} + \beta \begin{pmatrix}0 \\ 1\end{pmatrix}
= \begin{pmatrix}\alpha \\ \beta\end{pmatrix}
\]
\end{frame}

% --- Slide 11 ---
\begin{frame}{Kvantne logičke operacije}
\[
U = \begin{pmatrix} a & b \\ c & d \end{pmatrix}, \quad U^\dagger U = I
\]
\begin{itemize}
    \item Kvantne operacije su unitarne i reverzibilne.
    \item Za razliku od klasičnih logičkih kola (AND, OR), ne gube informaciju.
\end{itemize}
\end{frame}

% --- Slide 12 ---
\begin{frame}{Kvantna kola}
\begin{itemize}
    \item Sekvence unitarnih operacija nad kubitima.
    \item Primeri:
    \begin{itemize}
        \item Hadamardovo kolo
        \item CNOT
        \item Pauli-X, Y, Z
        \item Phase gate
        \item Toffoli
    \end{itemize}
\end{itemize}
\end{frame}

% --- Slide 13 ---
\begin{frame}[fragile]{Primer kvantnog kola u QisKit}
\begin{lstlisting}[language=Python]
import numpy as np
from qiskit import QuantumCircuit

qc = QuantumCircuit(3)
qc.h(0)
qc.p(np.pi / 2, 0)
qc.cx(0, 1)
qc.cx(0, 2)
print(qc.draw('text'))
\end{lstlisting}
\end{frame}

% --- Slide 14 ---
\begin{frame}{Hadamardovo kolo}
\[
H = \frac{1}{\sqrt{2}} \begin{pmatrix} 1 & 1 \\ 1 & -1 \end{pmatrix}
\]
\[
H|0\rangle = \frac{1}{\sqrt{2}}(|0\rangle + |1\rangle) = |+\rangle
\]
\[
H|1\rangle = \frac{1}{\sqrt{2}}(|0\rangle - |1\rangle) = |-\rangle
\]


Inverzna operacija: $H^{-1} = H$
\end{frame}

\begin{frame}
    \frametitle{Hadamardovo kolo sa n kubita}

    $$
H^{\otimes N} = \frac{1}{\sqrt{2^N}} \begin{pmatrix}
1 & 1 & 1 & \cdots & 1 \\
1 & -1 & 1 & \cdots & -1 \\
1 & 1 & -1 & \cdots & -1 \\
\vdots & \vdots & \vdots & \ddots & \vdots \\
1 & -1 & -1 & \cdots & 1
\end{pmatrix} $$


$$ H^{\otimes N}_{i,j} = \frac{(-1)^{i \cdot j}}{\sqrt{2^N}} $$
\end{frame}

\begin{frame}
    \frametitle{Superpozicija svih stanja sa n kubita}


    $$ H^{\otimes N} |0\rangle^{\otimes N} = \frac{1}{\sqrt{2^N}} \sum_{x=0}^{2^N-1} |x\rangle $$


\end{frame}

\begin{frame}
    \frametitle{Pauli-X kolo}


$$ X|0\rangle = |1\rangle $$

$$ X|1\rangle = |0\rangle $$


$$ X = \begin{pmatrix}
0 & 1 \\
1 & 0
\end{pmatrix} $$

$$ X^{\otimes N} = X \otimes X \otimes \cdots \otimes X $$

$$ X^{\otimes N} |x\rangle = |(2^N - 1) - x\rangle $$    

$$ X^{-1} = X $$

\end{frame}

\end{document}
